\chapter{Discusión}\label{cap:6}

\section{Síntesis e interpretación}\label{sec:sintesis}

Tras igualar la frecuencia de días húmedos (criterio i), ALADIN reproduce el gradiente norte--sur de persistencia de CR2MET, pero subestima la duración de los dry spells, sobre todo en la cola (t99). La calibración de R₀ corrige gran parte del sesgo asociado al exceso de días con precipitación débil, un problema persistente en modelos climáticos, que tienden a llover con demasiada frecuencia y poca intensidad (Sun et al., 2006; Stephens et al., 2010; Chen et al., 2021; Martinez-Villalobos et al., 2022). Sin embargo, no elimina las diferencias en duración. Igualar la frecuencia de días húmedos es, por tanto, necesario pero no suficiente: también importa cómo se organizan esos días en el tiempo y cuánto fragmentan las rachas secas. Este sesgo de cola es coherente con Petrova et al.~(2024), quienes encuentran que la mayoría de los modelos CMIP5 y CMIP6 subestiman la racha seca anual más larga observada y que los modelos que la subestiman en el histórico tienden a proyectar aumentos futuros menores. Aunque esa relación se obtuvo con modelos globales y un indicador anual, sugiere que el alargamiento proyectado por ALADIN podría ser conservador, especialmente en el norte árido, donde el acortamiento de la cola histórica es mayor.

Hacia 2040--2074 bajo SSP5-8.5, ALADIN proyecta menos dry spells, pero más largos: aumenta el número de rachas de 20 días o más, con señal significativa en Coquimbo y La Araucanía, y las PDFs de Maipo, Maule y Biobío extienden la cola. Esta reorganización es opuesta a la que Latoja et al.~(2025) documentan con CR2MET en Chile central (32°--40°S) entre 1960 y 2021: más dry spells, pero más cortos, junto con menos precipitación por racha húmeda. Ambos resultados describen un secado, pero no la misma reorganización temporal. En el registro observado, la disminución de precipitación se expresa sobre todo en eventos húmedos más débiles e intermitentes. En la proyección, se expresa en menos interrupciones de los períodos secos, coherente con el papel dominante del aumento de días secos en el secado proyectado para climas mediterráneos (Polade et al., 2014, 2017). Parte del contraste puede deberse al forzante: la megasequía combina variabilidad interna con una contribución antrópica (Boisier et al., 2016; Garreaud et al., 2020), mientras que bajo SSP5-8.5 domina la respuesta forzada. También difieren la métrica y la ventana: Latoja et al.~analizan la tendencia de la duración media en seis décadas, y aquí se compara la cola entre dos períodos de 35 años. Secarse, por tanto, no implica una única forma de reorganizar las rachas secas, y el cambio en su duración no puede inferirse solo del signo del cambio en la precipitación total.

\section{Implicaciones para evaluación de modelos y riesgo de sequía}\label{sec:implicaciones}

Los resultados históricos tienen una implicación directa para la evaluación de RCMs en climas mediterráneos y áridos. Validar la precipitación media o la frecuencia de días húmedos no basta para juzgar la persistencia seca: aun después de igualar esa frecuencia, ALADIN acorta las rachas y, sobre todo, su cola. Una evaluación orientada a sequía debería, por tanto, comparar la distribución de duraciones tras controlar la frecuencia de días húmedos (Leung et al., 2022) y reportar explícitamente la cola. Esto importa también para las proyecciones: si, como encuentran Petrova et al.~(2024), los modelos que subestiman las rachas largas históricas tienden a proyectar aumentos menores, el sesgo de cola histórico condiciona cuánto confiar en la magnitud del cambio futuro. La metodología, basada en PDFs que conservan el número de spells, risk ratios de conteos y calibración previa de días húmedos, es transferible a otros RCMs del dominio CHP.

Para la gestión hídrica y la agricultura, el resultado central es que el cambio proyectado se concentra en las rachas largas y no se deduce del cambio en la precipitación total. Un aumento en el número de rachas de 20 días o más puede tener consecuencias que la reducción de la precipitación anual no anticipa, porque los períodos prolongados sin lluvia condicionan el déficit de humedad del suelo, la recarga y la planificación del riego. Un índice como el CDD, que retiene un solo extremo por año, tampoco captura ese cambio, que se distribuye a lo largo de la cola.

\section{Limitaciones y perspectivas}\label{sec:limitaciones}

Varias limitaciones acotan la interpretación. Primera, se usa un único RCM y un único escenario de altas emisiones; los cambios proyectados describen la respuesta de ALADIN bajo SSP5-8.5, no la incertidumbre multimodelo. Segunda, el umbral R₀ calibrado en el histórico se aplica también al futuro, lo que supone que el sesgo de frecuencia de días húmedos del modelo no cambia con el clima. Tercera, los píxeles de un mismo pool no son independientes, de modo que el número de spells sobrestima la información independiente; por eso la incertidumbre se estima remuestreando años y no spells individuales. Cuarta, en el norte árido ALADIN acorta la cola respecto de CR2MET incluso tras calibrar R₀, de modo que la frecuencia absoluta de spells muy largos en Loa debe leerse como conservadora. Quinta, en regiones muy húmedas como Los Lagos, los RR para umbrales altos son inestables: los pocos eventos históricos producen intervalos de confianza amplios aunque el cociente puntual sea alto. Finalmente, el trabajo caracteriza la reorganización de las rachas secas, pero no diagnostica los mecanismos de circulación que la producen.

El trabajo futuro puede avanzar en tres direcciones. La primera es extender el análisis a otros RCMs del dominio CHP y a otros escenarios, lo que permitiría además evaluar si, como sugieren Petrova et al.~(2024), los modelos con mayor sesgo de cola histórico proyectan aumentos menores. La segunda es evaluar si ALADIN reproduce la tendencia observada hacia spells más numerosos y más cortos en Chile central (Latoja et al., 2025), y vincular el contraste con el futuro a cambios en la circulación. La tercera es relacionar las PDFs de duración con índices de sequía y caudales, para traducir el cambio en la cola a impactos hidrológicos.

Los resultados muestran que el secado futuro proyectado por ALADIN no consiste únicamente en una reducción de la precipitación acumulada. El modelo proyecta una reorganización temporal hacia menos interrupciones y rachas secas más prolongadas, con cambios particularmente visibles en la cola de la distribución en Chile central y centro-sur. Caracterizar esa persistencia permite identificar una dimensión del cambio hidroclimático que queda parcialmente oculta en la precipitación media y en índices basados en un único extremo anual.

\pendiente{ampliar con más literatura, posibles mecanismos de circulación, el papel
de los aerosoles en SSP3-7.0 y las implicancias para la gestión.}
